\section{Hidden-Block State Conversion}
\label{sec:state-conversion}

The central paths of \cref{sec:wta-central-path,sec:condensation-multiwinner}
concentrate trace on a hidden set of winning blocks.  This section studies the query cost of
preparing such a central state to trace-distance accuracy $\varepsilon$ when
the number of winners is promised to be either zero or $k$.  The two
extremes already appear in \cref{sec:condensation-exact-k-preparation}: a
sufficiently large error allows the public all-loser output, whereas a
sufficiently accurate coherent preparation costs as much as search.  In
between, the collision statistic in
\eqref{eq:cond-multi-collision} loses a factor $1/k$, and a bound on
pairwise trace distance gives no single measurement that works for every
unknown winner set.  Instead, we measure the label of the prepared state and
verify it with fresh queries.  This verified-sample decoder accepts a
zero-winner input only through verifier error and applies directly to
trace-close mixed outputs.

For parity blocks, the polynomial method then gives the query cost at every
nontrivial accuracy.  The key step is that averaging over inputs with fixed
block parities removes every character supported on a proper nonempty part
of a block, which divides the polynomial degree by the block length.  We
also compare the coupled centers with a batch of uncoupled qutrit-holonomy
SDPs.  Returning all component values admits no amortization, whereas
estimating their average has the usual approximate-counting query
complexity.

\subsection{The zero-versus-\texorpdfstring{$k$}{k} conversion problem}
\label{sec:state-conversion-model}

Let $f:\mathcal X_0\sqcup\mathcal X_1\to\{0,1\}$ be a possibly partial
Boolean predicate on one block.  Its bounded-error query complexity and
negative-weight adversary value satisfy
\[
 Q(f)=\Theta(A),\qquad A=\operatorname{Adv}^{\pm}(f).
\]
There are $G$ independently queried blocks with inputs $x_1,\ldots,x_G$; in
\cref{sec:winner-take-all}, block $g$ is group $g$, whose input is its sign
string $(\sigma_{g,1},\ldots,\sigma_{g,N})$, not one of its PSD blocks
$X_{g,i}$.  The winner set
$S=\{g:f(x_g)=1\}$ is promised to have size zero or $k$, where
$1\le k<G$.  One raw query addresses one coordinate of one block.  The
target on a $k$-winner input is
\begin{equation}
 \rho_S=\frac pk\sum_{g\in S}\ket g\!\bra g\otimes\rho_W
 +\frac{1-p}{G-k}\sum_{g\notin S}\ket g\!\bra g\otimes\rho_L,
 \qquad
 \rho_0=\frac1G\sum_{g=1}^G\ket g\!\bra g\otimes\rho_L^0,
 \label{eq:sc-targets}
\end{equation}
where $p=p_{G,k}$ is the winner mass and the constant-dimensional internal
states $\rho_W$, $\rho_L$, and $\rho_L^0$ are public.
Write
\begin{equation}
 \eta_{\rm int}=\Dtr(\rho_L,\rho_L^0),\qquad \delta=p-\varepsilon.
 \label{eq:sc-parameters}
\end{equation}
Here \(\eta_{\rm int}\) is the local internal-state discrepancy, not the
residual tolerance \(\eta\) used elsewhere.
The task is to output, on every promised input, a state within trace distance
$\varepsilon$ of its target.  The upper bounds below use a clean coherent
evaluator as in \eqref{eq:cond-clean-evaluator}, which computes the
membership bit for $S$, at cost $T_{\rm eval}$ per call.  An exact evaluator
costs $\Theta(Q_E(f))$ raw queries, where $Q_E$ denotes exact quantum query
complexity.  A bounded-error evaluator may be used instead after coherent
error reduction to operator-norm error below the per-call budget of the
amplification, at cost
$O(Q(f)\log(G/\min\{\delta,\varepsilon\}))$ per call; the tildes absorb
this factor.  For $f=\operatorname{PARITY}_N$, $Q_E(f)=\lceil N/2\rceil$, and
$Q(f)$ and $\operatorname{Adv}^{\pm}(f)$ are both $\Theta(N)$.  As in the
condensation applications, we assume $p\ge k/G$ when invoking rejection
preparation.
For classical algorithms we count only queries; preparing the output state
from the classical transcript is free.

Consider the test that measures the label of an $\varepsilon$-accurate
output and then evaluates $f$ on the sampled block.  On a $k$-winner input,
it accepts with probability at least $p-\varepsilon=\delta$, ignoring
verifier error.  On a zero-winner input, every proposed block is a loser, so
only verifier error can cause acceptance, whatever state the circuit
produces.  Unlike the collision statistic, this test is a single
measurement that works for every winner set $S$.

\begin{theorem}[Verified-sample lower bound]
\label{thm:sc-m1}
Fix $\delta\in(0,1)$.  Suppose a $T$-query circuit solves
\eqref{eq:sc-targets} with $\varepsilon\le p-\delta$.  If
\begin{equation}
 \sqrt{G/k}\ge C\delta^{-1}\log(e/\delta)
 \label{eq:sc-m1-domain}
\end{equation}
for a sufficiently large universal constant $C$, then
\[
 T=\Omega\!\left(\delta A\sqrt{G/k}\right).
\]
For $f=\operatorname{PARITY}_N$, this is
$T=\Omega(\delta N\sqrt{G/k})$.
\end{theorem}

\begin{proof}
First reduce the zero-versus-one problem of \cref{thm:cond-unique-or} to
zero-versus-$k$.  Put $G'=\lfloor G/k\rfloor$.  Replace each of the $G'$
input blocks by $k$ addressable copies and pin the remaining
$G-kG'<k$ blocks to a fixed public element of $\mathcal X_0$.  A query to a
copy is simulated by one query to its original.  Zero winners remain zero,
and a unique winner produces exactly $k$ winners.  The Cartesian-product
promise is preserved.  Since \eqref{eq:sc-m1-domain} forces $G/k$ to be at least a large
constant, $G'=\Theta(G/k)$, and
\eqref{eq:cond-unique-or-adv} gives a decision lower bound
$\Omega(A\sqrt{G/k})$.  This construction also works when $k$ does not divide $G$.

Now repeat the following trial $R=\Theta(1/\delta)$ times.  Run the
preparation circuit afresh, measure its label $g$, and evaluate $f(x_g)$
with error at most $\beta=\delta/8$.  Standard error reduction costs
$O(Q(f)\log(e/\delta))$ queries per verification.  On a $k$-winner input,
contractivity of trace distance under the two-outcome measurement
$\{\sum_{g\in S}\ket g\!\bra g\otimes I,I-\cdots\}$ gives
\[
 \Pr[g\in S]\ge p-\varepsilon\ge\delta,
 \qquad \Pr[\text{accept}]\ge\delta(1-\beta)\ge\delta/2.
\]
On a zero-winner input, every label fails the true predicate, regardless of
the prepared state, and $\Pr[\text{accept}]\le\beta=\delta/8$.  A
multiplicative Chernoff bound distinguishes these two rates with constant
error using $R=\Theta(1/\delta)$ trials.  Its total cost
\[
 O\!\left(\delta^{-1}
   [T+Q(f)\log(e/\delta)]\right)
\]
is therefore $\Omega(A\sqrt{G/k})$.  Since $Q(f)=\Theta(A)$,
\eqref{eq:sc-m1-domain} absorbs the verification term and proves the claim.
\end{proof}

The domain condition \eqref{eq:sc-m1-domain} is restrictive.  At
constant $\delta$ it requires only $G/k=\Omega(1)$; \cref{cor:sc-m3}
assumes the stronger uniform condition $k=O(G/\log^2G)$.  In general it
requires $G/k=\Omega(\delta^{-2})$ up to logarithms.  In exchange,
\cref{thm:sc-m1} holds for a general inner predicate $f$ and assumes no
access to a purification or to the inverse of the preparation.

\begin{theorem}[Detuned rejection preparation]
\label{thm:sc-m2}
Assume $p\ge k/G$, $1-p\ge c_0>0$, and $0<\varepsilon<p$.  With
$\delta=p-\varepsilon$, if $4\eta_{\rm int}\le\varepsilon$, then
\begin{equation}
 Q_\varepsilon=O\!\left(
 T_{\rm eval}\left[1+\sqrt{\delta G/k}\right]c_0^{-1/2}
 \log\frac{e}{\min\{\delta,\varepsilon\}}\right).
 \label{eq:sc-m2-upper}
\end{equation}
If $0<\varepsilon<4\eta_{\rm int}$, then
\begin{equation}
 Q_\varepsilon=O\!\left(
 T_{\rm eval}\sqrt{G/k}\log(e/\varepsilon)\right).
 \label{eq:sc-m2-fallback}
\end{equation}
When $\varepsilon=\Omega(\delta)$, the logarithm
in \eqref{eq:sc-m2-upper} is simply $O(\log(e/\delta))$.
\end{theorem}

\begin{proof}
We prepare with a detuned winner mass $m$, chosen in two regimes:
\begin{equation}
 \begin{aligned}
 m&=\begin{cases}
 \max\{3\delta/2,k/G\},&\delta\le2p/3,\\
 \max\{p-\varepsilon/2,k/G\},&\delta>2p/3,
 \end{cases}\\[-2pt]
 m&\le p:\quad
 \begin{cases}
 3\delta/2\le p\ \text{and}\ k/G\le p,&\delta\le2p/3,\\
 p-\varepsilon/2\le p\ \text{and}\ k/G\le p,&\delta>2p/3.
 \end{cases}
 \end{aligned}
 \label{eq:sc-detuning}
\end{equation}
Then $k/G\le m\le p$, $m-\delta\ge\delta/2$ in the first
regime, and $m-\delta\ge\varepsilon/2$ in the second.  Moreover
$\sqrt{Gm/k}=O(1+\sqrt{\delta G/k})$: this is immediate in the first
regime, while $m=\Theta(\delta)$ in the second.  The cap $m\le p$
is essential: with $1-p\ge c_0$, it keeps the zero-winner acceptance
probability from adding a factor $(1-m)^{-1/2}$ to the amplification cost.

Run the rejection construction of \cref{prop:cond-rejection-preparation}
with the label copied before amplification and the copy included in every
reflection, and with detuned weights
\[
 a'_W=m/k,\qquad a'_L=(1-m)/(G-k),\qquad a'_{\max}=m/k.
\]
Starting from the uniform label superposition, copy the label to an
environment register, compute the membership bit, rotate an acceptance
qubit with amplitude $\sqrt{a'_b/a'_{\max}}$, and conditionally attach
public purifications of $\rho_W$ and $\rho_L$.  A trial has acceptance
probability
\begin{equation}
 P_k=\frac{k}{Gm},\qquad
 P_0=\frac{k(1-m)}{m(G-k)}
 \label{eq:sc-branch-acceptance}
\end{equation}
on $k$-winner and zero-winner inputs, respectively.  Both are at least a
constant times $c_0k/(Gm)$, because $m\le p$ implies
$1-m\ge1-p\ge c_0$.

Apply fixed-point amplification simultaneously to the two possible
acceptance angles \cite{yoder2014-fixed-point-quantum-search}.  The
fixed-point reflections act on every register the trial touches, including
the environment copy of the label, which is uncomputed and recomputed inside
each reflection.  With
\[
 r=\frac18\min\{\delta,\varepsilon\},
\]
the acceptance-subspace residual can be made at most $r^2$ on both kinds of input
using
\[
 O\!\left(\sqrt{Gm/k}\,c_0^{-1/2}\log(e/r)\right)
\]
evaluator calls.  The quadratic residual is necessary: overlap error of a
purification becomes its square root in trace distance.  Writing $q=k/G$,
\eqref{eq:sc-branch-acceptance} gives
$P_0/P_k=(1-m)/(1-q)$.  Thus $m=q$ gives equal acceptance probabilities.  When
$q=o(m)$, however, an exact schedule tuned to the larger probability $P_k$
under-rotates a zero-winner input and can leave $\Theta(m^2)$ junk probability,
which exceeds the error budget for constant $p$.  Fixed-point amplification avoids
this mismatch.

Conditioned on acceptance, the output on a $k$-winner input is exactly
\eqref{eq:sc-targets} with winner mass $m$ in place of $p$; the zero-winner
output is exactly $G^{-1}\sum_g\ket g\!\bra g\otimes\rho_L$.  Hence the
respective errors are at most
\[
 (p-m)+r,\qquad \eta_{\rm int}+r.
\]
If $\delta\le2p/3$, then $\varepsilon=p-\delta\ge\delta/2$ and
$p-m\le\varepsilon-\delta/2$, so both errors are below
$\varepsilon$.  If $\delta>2p/3$, then $p-m\le\varepsilon/2$ and
$r\le\varepsilon/8$, so again both errors are at most $\varepsilon$.  Substituting the bound on
$m$ gives \eqref{eq:sc-m2-upper}.

For \eqref{eq:sc-m2-fallback}, first decide whether there are zero or $k$
winners by fixed-point search, reducing its error to $O(\varepsilon)$.  If
the decision is $k$, use the exactly-$k$ circuit \eqref{eq:cond-ap3}; if it
is zero, prepare $\rho_0$ directly.  The decision error contributes at most
$O(\varepsilon)$ in trace distance, and the stated cost dominates both
steps.  Deciding first avoids circular use of the exactly-$k$ promise.
\end{proof}

The precision logarithm in \eqref{eq:sc-m2-upper} depends on
$\min\{\delta,\varepsilon\}$, not on $\delta$ alone.  In the second regime of \eqref{eq:sc-detuning}, a
fixed-point residual of order $\delta$ would not suffice when
$\varepsilon\ll\delta$; the residual must be $O(\varepsilon)$.

\begin{corollary}[Constant-visibility regime]
\label{cor:sc-m3}
Let $f=\operatorname{PARITY}_N$, let $p$ be bounded away from zero and one,
and suppose
\[
 k=O(G/\log^2G),\qquad
 4\eta_{\rm int}\le\varepsilon\le(1-c)p
\]
for a constant $c>0$.  Then, up to fixed-point precision logarithms,
\[
 Q_\varepsilon=\Thetatilde\!\left(N\sqrt{G/k}\right).
\]
By \eqref{eq:sc-m2-fallback}, the same scale holds for
$\varepsilon<4\eta_{\rm int}$ when $\delta=\Omega(1)$.  If $\varepsilon$
exceeds the minimax radius of the target family, zero queries suffice; when
the internal loser states agree, \eqref{eq:cond-multi-visibility} bounds
the public output's error by $\max\{p,k/G\}$.
\end{corollary}

\begin{proof}
Here $\delta=p-\varepsilon=\Omega(1)$ and $A=\Theta(N)$.
The domain assumption makes \eqref{eq:sc-m1-domain} valid, so
\cref{thm:sc-m1} gives the lower bound.  The upper bound is
\cref{thm:sc-m2}.  The zero-query statement follows by returning the public
all-loser target and using \eqref{eq:cond-multi-visibility}.
\end{proof}

\begin{theorem}[Classical query complexity]
\label{thm:sc-m4}
For $f=\operatorname{PARITY}_N$, suppose $p\ge k/G$,
$1-p=\Omega(1)$, and
\begin{equation}
 2\eta_{\rm int}\le\varepsilon,
 \qquad Ck/G\le\delta=p-\varepsilon\le1/2.
 \label{eq:sc-m4-domain}
\end{equation}
Here $C>0$ is a sufficiently large universal constant.
Then the randomized classical raw-query complexity is
\[
 R_\varepsilon=\Theta(N\delta G/k).
\]
\end{theorem}

\begin{proof}
For the upper bound, inspect
$t=\Theta(\delta G/k)$ uniformly random blocks without replacement and
compute each parity exactly using $N$ queries.  Put $q=k/G$.  The probability
of finding at least one winner is
\[
 P_{\rm find}=1-\frac{\binom{G-k}{t}}{\binom Gt}.
\]
The constant in $t$ can be chosen so that
$P_{\rm find}\ge\delta$ throughout
\eqref{eq:sc-m4-domain}.  This use of a linear budget relies on
$\delta\le1/2$; as $\delta\to1$, the correct sampling budget contains
$(G/k)\log(1/(1-\delta))$.

Let $E_{\rm all}$ be the event that every sampled block is a winner.  Its
public probability is at most $q^t\le q\le\delta$.  On $E_{\rm all}$ output
a uniformly selected sampled winner.  On
$E_{\rm find}\setminus E_{\rm all}$, do the same with a public probability
chosen so that the total probability of winner output is exactly $\delta$;
this is possible because
$\Pr(E_{\rm all})\le\delta\le\Pr(E_{\rm find})$.  In all other cases select
uniformly among sampled losers.  A loser is available whenever the latter
rule is used.  Exchangeability makes the selected winner uniform on $S$ and
the selected loser uniform on its complement.  Attaching the corresponding
public internal state therefore prepares exactly the target shape with
winner mass $\delta$ in place of $p$.  Its trace error on a $k$-winner input is
$p-\delta=\varepsilon$.  On a zero-winner input the output is the uniform-label
$\rho_L$ state and has error $\eta_{\rm int}\le\varepsilon/2$.  Thus
$O(N\delta G/k)$ queries suffice.

For the lower bound, append exact $N$-query parity verification to a sampled
output label and repeat the preparation $\Theta(1/\delta)$ times.  This
decides zero versus $k$ winners with one-sided error and cost
$O(\delta^{-1}(T+N))$.  For the decision lower bound, use the $k$-copy
embedding from \cref{thm:sc-m1} and draw each original block uniformly
conditioned on its parity.  Every proper subset of a block is uniform, so a
classical transcript is independent of the planted parity until all $N$
positions of that block have been queried.  A $q$-query tree completes at
most $q/N$ original blocks, and it hits the uniformly planted winner with
probability at most $O(qk/(NG))$.  Yao's principle therefore gives
$\Omega(NG/k)$ decision queries.  Consequently
\[
 \delta^{-1}(T+N)=\Omega(NG/k).
\]
The lower restriction $\delta\ge Ck/G$ absorbs the additive $N$ term.
\end{proof}

We claim no lower bound when $\delta=o(k/G)$: there, a public output with
uniform label may already meet the accuracy requirement.  The upper
restriction $\delta\le1/2$ in \eqref{eq:sc-m4-domain} is also needed:
near one, the probability of finding a winner is no longer linear in the
number of sampled blocks.

\subsection{Composed small-success search}
\label{sec:state-conversion-small-success}

The verified-sample argument loses a factor $\sqrt\delta$ by repeating the
preparation.  For parity blocks, the polynomial method recovers it.  We
state the degree-division step as a lemma because
\cref{sec:state-conversion-batch} also uses it.

\begin{lemma}[Parity collapse]
\label{lem:sc-parity-collapse}
Let a $T$-query quantum algorithm access
$x\in\{\pm1\}^{G\times N}$, and let $P(x)$ be the probability of a fixed
final measurement event.  For block parities
$z_g=\prod_{j=1}^N x_{g,j}$, define
\[
 q(z)=\mathbb E[P(x)\mid \textstyle\prod_jx_{g,j}=z_g
                  \text{ for every }g].
\]
Here the conditional expectation is uniform over each parity fiber,
independently across blocks.
Then $q$ is a real multilinear polynomial in $z\in\{\pm1\}^G$ of degree
at most $\lfloor2T/N\rfloor$.  If $0\le P\le1$, then $0\le q\le1$ on
the Boolean cube.
\end{lemma}

\begin{proof}
The polynomial method writes $P$ as a real multilinear polynomial of degree
at most $2T$ \cite{beals2001-quantum-lower-bounds-by-polynomials}.  Consider
a monomial $x_R$.  Conditional expectation factors over blocks.  In one
block, the average of a character on a parity coset is one when its share of
$R$ is empty, is $z_g$ when that share is the full block, and is zero for
every proper nonempty share.  The last assertion follows by pairing inputs
that differ in two coordinates, one inside and one outside the proper
share.  Thus a raw monomial with nonzero average becomes a monomial in $z$ whose
degree is the number of full blocks it contains, at most
$|R|/N\le2T/N$.  Averaging all monomials proves the claim.
\end{proof}

\begin{remark}[Computational check]
We checked the character calculation in \cref{lem:sc-parity-collapse} by
exhaustive enumeration for $N=2,3,4,5$, and the full identity on random
multilinear polynomials.
\end{remark}

\begin{theorem}[One-sided composed small-success bound]
\label{thm:sc-small-success}
Suppose a $T$-query quantum algorithm accepts with probability at most
$\beta\le\delta/2$ on every all-even input and at least $\delta$ on every
input with exactly $k$ odd-parity blocks.  Then, for absolute constants
$c_1,c_2>0$,
\begin{equation}
 T\ge c_1N\sqrt{\delta G/k}-c_2N.
 \label{eq:sc-small-success}
\end{equation}
\end{theorem}

\begin{proof}
Let $P(x)$ be the acceptance probability.  By
\cref{lem:sc-parity-collapse}, conditioning on all block parities gives a
polynomial $q(z)$ of degree at most
$D=\lfloor2T/N\rfloor$, bounded between zero and one at every Boolean
input.  It is at most $\beta$ at the all-even point and at least $\delta$
at every parity vector of odd-block weight $k$.

Average $q$ over all block permutations.  Minsky--Papert symmetrization
\cite{minsky1988-perceptrons} gives a univariate polynomial $h$ of degree at
most $D$ such that
\[
 0\le h(w)\le1\quad(w=0,1,\ldots,G),\qquad
 h(0)\le\beta,\qquad h(k)\ge\delta.
\]
Fix a sufficiently small absolute $\gamma>0$.  First suppose
$D\le\gamma\sqrt G$.  The Coppersmith--Rivlin inequality bounds a
degree-$D$ polynomial bounded at the $G+1$ equally spaced integer points by
\[
 \sup_{t\in[0,G]}|h(t)|\le C(\gamma),
\]
where $C(\gamma)$ is an absolute constant in this degree range
\cite{coppersmith1992-growth-polynomials,klauck2007-quantum-classical-sdpt}.
We use only the existence of the universal Coppersmith--Rivlin constants,
which the original paper does not give numerically; this is the form quoted
in Theorem~8 of Klauck--\v Spalek--de Wolf.  Erd\'elyi proves a related
explicit-constant inequality in a different, single-point-versus-supremum
parametrization
\cite{erdelyi2016-coppersmith-rivlin}.

Markov's inequality, rescaled from $[-1,1]$ to $[0,G]$, now gives
\[
 \sup_{[0,G]}|h'|\le\frac{2D^2C(\gamma)}G.
\]
The mean value theorem on $[0,k]$ gives a point $\xi$ with
\[
 h'(\xi)\ge\frac{h(k)-h(0)}k
 \ge\frac{\delta-\beta}k\ge\frac{\delta}{2k}.
\]
Therefore $D=\Omega(\sqrt{\delta G/k})$.  In the complementary case
$D>\gamma\sqrt G$, the same conclusion follows directly from
$\delta\le1\le k$.  Finally, $D\le2T/N$, with an additive $O(N)$ allowance for integer
rounding, proves
\eqref{eq:sc-small-success}.
\end{proof}

To our knowledge, \cref{thm:sc-small-success} is new as stated, but its
ingredients are standard and experts could plausibly derive it.  It is not
a new general composition theorem.

\begin{corollary}[Composed search]
\label{cor:sc-s1}
If a quantum algorithm outputs an odd-parity block with probability at least
$\delta$ whenever exactly $k$ blocks are odd, then
\[
 T=\Omega\!\left(\sqrt\delta\,N\sqrt{G/k}\right)-O(N).
\]
This is tight up to constants and the additive $N$ term.
\end{corollary}

\begin{proof}
Verify the output block by computing its parity exactly with $N$ fresh
queries.  Acceptance is then zero on an all-even input and at least
$\delta$ on an exactly-$k$ input, so \cref{thm:sc-small-success} applies to
$T+N$ queries.  Conversely, if $\delta\le4k/G$ a uniformly random label is
odd with probability $k/G\ge\delta/4$ at zero queries.  Otherwise $k/G<1/4$,
so the initial Grover angle $\theta=\arcsin\sqrt{k/G}$ is below $\pi/6$, and
$t=\lceil c\sqrt{\delta G/k}\rceil\ge1$ Grover iterations with $c=1/6$ keep
$(2t+1)\theta\le\pi/2$; each uses a $\Theta(N)$-query parity marker, and
together they produce a winner with probability
$\sin^2((2t+1)\theta)\ge\sin^2(2c\sqrt\delta)=\Theta(\delta)$.  Direct exact
verification costs another $N$ queries.
For one marked item, Zalka proved exact optimality
\cite{zalka1999-grovers-quantum-searching-algorithm-is}.  The small-success
scaling for $k$ marked items follows from the angle bound in Theorem~2 of
the published Dohotaru--H\o yer paper (Theorem~9 of the arXiv version) plus
the standard $k$-copy embedding.
Carolan--Poremba Corollary~2.4 records the packaged inequality
$\delta\le8(T+1)^2k/G$; our attribution uses the angle theorem and embedding,
not their intermediate distance-based derivation
\cite{dohotaru2009-exact-grover,carolan2024-quantum-one-wayness}.  Zalka's
result alone does not establish the $k$-marked statement.  Aaronson's
direct-product Theorem~4.6 instead concerns the find-all-$k$ analogue
\cite{aaronson2005-limitations-quantum-advice}.
\end{proof}

\begin{corollary}[Full conversion curve]
\label{cor:sc-s2}
Let $f=\operatorname{PARITY}_N$, $p\ge k/G$, $1-p=\Omega(1)$,
$4\eta_{\rm int}\le\varepsilon<p$, and $\delta=p-\varepsilon\ge Ck/G$ for a
sufficiently large constant $C$.  Then
\begin{equation}
 \boxed{Q_\varepsilon=\Thetatilde\!\left(
 N\sqrt{(p-\varepsilon)G/k}\right).}
 \label{eq:sc-full-curve}
\end{equation}
The tilde hides only fixed-point precision logarithms: the hidden factor is
$\log(e/\min\{\delta,\varepsilon\})=\log(1/\varepsilon)+O(1)$ when
$\varepsilon\ll\delta$, and can be unbounded relative to the main term for
superpolynomially small $\varepsilon$.  The lower bound is nonvacuous once
$\delta G/k$ is a sufficiently large constant.
\end{corollary}

\begin{proof}
One preparation followed by label measurement and exact parity verification
uses $T+N$ queries.  It accepts with probability zero on a zero-winner input
and at least $p-\varepsilon=\delta$ on a $k$-winner input.  Applying
\cref{thm:sc-small-success} and absorbing its additive $O(N)$ term under the
domain assumption gives the lower bound.  The upper bound is
\cref{thm:sc-m2}; its additive one is absorbed on the same domain.
\end{proof}

Unlike \cref{thm:sc-m1}, this conclusion needs neither
$k=O(G/\log^2G)$ nor $G/k=\widetilde\Omega(\delta^{-2})$.  It is
parity-specific: the gain comes from \cref{lem:sc-parity-collapse}.

\begin{corollary}[Mixed-output preparation profile]
\label{cor:sc-s3}
Suppose a $T$-query circuit, on every exactly-$k$ parity input, outputs a
state within trace distance $\varepsilon\le p/2$ of the winner-mass-$p$
target.  Then
\begin{equation}
 T=\Omega\!\left(\sqrt p\,N\sqrt{G/k}\right)-O(N).
 \label{eq:sc-mixed-profile}
\end{equation}
Thus the first two lines of the preparation profile \eqref{eq:cond-ap6}
are tight even for trace-close mixed outputs: the lower bounds are
$\Omega(N\sqrt{G/k})$ in the subcritical phase and
$\Omega(N(G/k)^{1/4})$ at criticality.  In the supercritical phase the
displayed lower bound is vacuous, so it does not show that the $O(N)$ upper
bound is tight.
\end{corollary}

\begin{proof}
Run the circuit also on an all-even input, which is outside its preparation
promise but on which the physical circuit remains defined.  Measure a label
and verify its parity exactly.  Acceptance is zero on this input and is at
least $p-\varepsilon\ge p/2$ on every exactly-$k$ input.  Apply
\cref{thm:sc-small-success} with $\delta=p/2$.  The three phase statements
follow by substituting the winner masses from
\cref{thm:cond-multiwinner}: $p=\Theta(1)$,
$p=\Theta(\sqrt{k/G})$, and $p=\Theta(k/G)$.
\end{proof}

\Cref{cor:sc-s3} weakens the hypothesis of the preparation lower bounds in
\cref{sec:condensation-exact-k-preparation}.  The lower bound \eqref{eq:cond-ap5} was obtained by coherent
extraction through \eqref{eq:cond-ap4}, and therefore assumed access to a
preparation unitary and its inverse.  For parity blocks and
$\varepsilon\le p/2$, \cref{cor:sc-s3} replaces this access by the weaker
requirement that the circuit output a trace-close mixed state.  The upper bound \eqref{eq:cond-ap3} and the caveat on the
supercritical phase in \eqref{eq:cond-ap6} are unchanged.  The family
of \cref{sec:wta-central-state} has parity blocks and $k=1$, so this
shows which of its condensed central states reveal the hidden holonomy; the
search-then-crossover approach \eqref{eq:wta-architecture} is unchanged.

Combining \cref{cor:sc-s2,thm:sc-m4} gives, on their common domain,
\begin{center}
\begin{tabular}{@{}lcc@{}}
\toprule
model & query complexity & dependence on $\delta=p-\varepsilon$ \\ \midrule
quantum & $\Thetatilde(N\sqrt{\delta G/k})$ & square root \\
randomized classical & $\Theta(N\delta G/k)$ & linear \\
\bottomrule
\end{tabular}
\end{center}
Thus at every nontrivial accuracy, from $\delta=\Theta(k/G)$ to
constant $\delta$, the classical-to-quantum ratio is
$\widetilde\Theta(\sqrt{\delta G/k})$, as for search.  It is
$\widetilde\Theta(1)$ near the public-output boundary $\delta=\Theta(k/G)$
and $\widetilde\Theta(\sqrt{G/k})$ at constant $\delta$.

\begin{openproblem}[General-inner composition]
\label{open:sc-general-composition}
Does every predicate $f$ satisfy a composed small-success lower bound
\[
 \Omega\!\left(\sqrt\delta\,Q(f)\sqrt{G/k}\right)
\]
for finding a winner with probability $\delta$?  The adversary reduction in
\cref{thm:sc-m1} gives only linear dependence on $\delta$, while the parity
proof uses the fact that every character on a proper nonempty part of a
block averages to zero.
\end{openproblem}

\subsection{Uncoupled holonomy batches}
\label{sec:state-conversion-batch}

The global trace row in \cref{sec:wta-construction} couples all components.
For contrast, take $G$ disjoint copies of the qutrit-holonomy component,
each with its own public trace-one row, and no equality that joins two
components.  The path length and local costs are those of
\cref{eq:wta-parameters,eq:wta-local-costs}.  Consequently the product has
$L=33GN$ constant-size PSD blocks, scalar row sparsity at most three, and
scalar column sparsity at most two.  If
\[
 h_g=\prod_{j=1}^N\sigma_{g,j},\qquad y_g=(1+h_g)/2,
\]
then \eqref{eq:wta-cost-spectra} gives component optima
$v_+=29/10$ and $v_-=14/5$.

\begin{theorem}[No amortization for vector output]
\label{thm:sc-batch-vector}
Returning estimates $\widehat v_1,\ldots,\widehat v_G$ satisfying
$|\widehat v_g-v_{h_g}|\le1/25$ for every $g$, with worst-case joint success
probability at least $2/3$, has quantum raw-query complexity
\[
 \Theta(GN)=\Theta(L).
\]
Moreover, there are constants $a,b>0$ such that every algorithm with at
most $aGN$ queries has worst-case joint success probability at most
$2^{-bG}$.
\end{theorem}

\begin{proof}
Threshold every estimate at $57/20$ to recover all $h_g$.  Their product is
the parity of all $GN$ raw signs, which has quantum query complexity
$\Omega(GN)$; reading all signs is a matching upper bound.  For the stronger
claim, specialize Lee--Roland's strong direct-product theorem to
$G$ copies of $\operatorname{PARITY}_N$
\cite{lee2012-a-strong-direct-product-theorem}.  Its Theorem~1.1 with
success parameter $3/4$ gives
\[
 Q_{1-(3/4)^{G/2}}(f^{(G)})
 \ge\frac{G\ln(9/8)}{8000}Q_{1/4}(f).
\]
Since $Q_{1/4}(\operatorname{PARITY}_N)=\Theta(N)$, with fewer than a
fixed constant times $GN$ queries, simultaneous decoding succeeds with
worst-case probability at most $(3/4)^{G/2}=2^{-\Omega(G)}$.
\end{proof}

Now average the $G$ component objectives into one SDP.  Since the components
share no constraint, its optimal value is
\begin{equation}
 V(\sigma)=\frac1G\sum_{g=1}^Gv_{h_g}
 =\frac{14}{5}+\frac1{10G}\sum_{g=1}^Gy_g.
 \label{eq:sc-batch-value}
\end{equation}

\begin{theorem}[Query complexity of the averaged value]
\label{thm:sc-batch-scalar}
For $0<\varepsilon\le1/100$,
\begin{align}
 Q_\varepsilon(V)&=\Theta\!\left(
 N\min\{G,1/\varepsilon\}\right),
 \label{eq:sc-batch-quantum}\\
 R_\varepsilon(V)&=\Theta\!\left(
 N\min\{G,1/\varepsilon^2\}\right).
 \label{eq:sc-batch-classical}
\end{align}
\end{theorem}

\begin{proof}
An exact coherent query to $y_g$ costs $\Theta(N)$ raw sign queries.
Amplitude estimation estimates its mean to error $10\varepsilon$ using
$O(1/\varepsilon)$ such calls
\cite{brassard2002-quantum-amplitude-amplification-and-estimation}; when
$1/\varepsilon>G$, compute all parities.  This proves the quantum upper
bound.  Classically, sample
$O(\min\{G,1/\varepsilon^2\})$ groups without replacement, compute their
parities, and use the empirical mean, scanning all groups when the minimum
is $G$.

For every $G$, restricting to inputs in which all $G$ groups carry the same
$N$-bit string $x$ makes $V$ change by $1/10>2\varepsilon$ when the parity of
$x$ flips, so an $\varepsilon$-estimator computes
$\operatorname{PARITY}_N$ and $T\ge N/2$; this settles $G<6$, where
$\min\{G,1/\varepsilon\}=O(1)$.  For $G\ge6$, choose Hamming weights
$\ell,\ell'\in[G/4,3G/4]$ separated by
$\Delta=\max\{1,\lceil21\varepsilon G\rceil\}$.  Their values in
\eqref{eq:sc-batch-value} differ by more than $2\varepsilon$, so thresholding
an estimator gives a bounded-error polynomial distinguishing the two
weights.  Nayak--Wu's symmetric approximate-counting lower bound gives
degree
\[
 \Omega\!\left(\frac{\sqrt{m(G-m)}}{\Delta}\right)
 =\Omega(\min\{G,1/\varepsilon\}),
\]
where $m$ is the selected weight farther from $G/2$
\cite{nayak1999-the-quantum-query-complexity-of}.  Applying the already
proved degree division in \cref{lem:sc-parity-collapse} forces
$\lfloor2T/N\rfloor=\Omega(\min\{G,1/\varepsilon\})$.

The same lower bound also follows from perfect negative-adversary
composition: compose the partial symmetric outer gap-counting problem with
$G$ independent copies of $\operatorname{PARITY}_N$ and combine Nayak--Wu's outer bound with the
Belovs--Lee composition theorem
\cite{belovs2020-the-quantum-query-complexity-of}.  Thus the parity-collapse
argument here is a construction-specific proof, not a new composition
theorem.

For the classical bound, we show that even an adaptive algorithm must read
every sign of a group to learn its parity.  Fix $h\in\{\pm1\}^G$ and draw the raw
signs of each group uniformly conditional on their product being $h_g$.
While at least two signs of a group remain unseen, a simulator answers
each new query to that group with an independent fair sign.  Only when the
final unseen sign in a group is queried does the simulator query $h_g$ and
force the required product.
This exactly reproduces the conditional distribution under arbitrary
interleaving.  Each parity-oracle query consumes $N$ distinct raw queries,
so a depth-$q$ raw tree induces a parity-bit tree of depth at most
$\lfloor q/N\rfloor$.  The classical bounded-error complexity of estimating
a $G$-bit mean to error $10\varepsilon$ is
$\Theta(\min\{G,1/\varepsilon^2\})$, by the standard Yao distribution on
two nearby central Hamming weights.  Multiplication by $N$ proves
\eqref{eq:sc-batch-classical}.
\end{proof}

The two bounds give three regimes:
\begin{center}
\begin{tabular}{@{}lcc@{}}
\toprule
accuracy & $Q_\varepsilon(V)$ & $R_\varepsilon(V)$ \\ \midrule
$\varepsilon=\Omega(G^{-1/2})$
 & $\Theta(N/\varepsilon)$ & $\Theta(N/\varepsilon^2)$ \\
$G^{-1}=O(\varepsilon)=O(G^{-1/2})$
 & $\Theta(N/\varepsilon)$ & $\Theta(NG)$ \\
$\varepsilon=O(G^{-1})$
 & $\Theta(NG)$ & $\Theta(NG)$ \\
\bottomrule
\end{tabular}
\end{center}
At $\varepsilon=\Theta(G^{-1/2})$, the costs are
$\Theta(N\sqrt G)$ and $\Theta(NG)$, exactly the quantum and classical
scales of the coupled winner-take-all value bound
\eqref{eq:wta-value-complexity}.  The quantum advantage therefore depends on how the accuracy scales with
$G$; it is not an exponential advantage.

\begin{corollary}[Classical output of all Newton directions]
\label{cor:sc-parallel-newton}
At the public point $X_{g,i}=I/3$ and path parameter
$\mu=1/(GP)$, the uncoupled batch has reduced root Hessian $(9/G)I$ and
condition number one.  Nevertheless, producing classical descriptions of
all $G$ root Newton directions to entrywise error less than $1/200$ costs
$\Theta(GN)$ queries, and query budgets below a suitable constant times $GN$ have
$2^{-\Omega(G)}$ worst-case joint success probability.
\end{corollary}

\begin{proof}
Using $A_h$ from \eqref{eq:wta-Ah}, direct elimination gives
\[
 \Delta Z_g=-\frac1{90}(H_{12}+H_{23}+h_gH_{13}).
\]
Every full physical step remains positive definite, with minimum eigenvalue
at least $14/45$.  The hidden $13$ entry has magnitude $1/90>1/200$, so its
sign reveals $h_g$.  Apply \cref{thm:sc-batch-vector}.
\end{proof}

These batch results concern raw queries to the signed copy rows.  An
oracle that supplies group parities, aggregate gradients, or component
optima has already done the hard part: combining the raw signs into
parities or aggregates of them.
\Cref{thm:sc-batch-vector} concerns writing out all components classically,
not preparing one normalized quantum state of the global direction.  Such a
state gives each group amplitude $G^{-1/2}$, so a preparation with constant
error need not reveal any single entry, as \cref{rem:wta-dilution} shows.  Finally, the claims assume exact
equalities; the path operator has small singular values, so they imply
nothing about solutions with constant residuals.

\subsection{Open questions}
\label{sec:state-conversion-open}

Three questions remain open.  First, the decoder based on winner mass does
not detect the internal-state discrepancy $\eta_{\rm int}$.  The minimax
radius and query complexity of the conversion problem when $\eta_{\rm int}$
matters, especially when $p$ is below the condensation threshold, are
open.  Second,
\cref{open:sc-general-composition} asks for the $\sqrt\delta$ composition
bound beyond parity.  Third, our tight bounds exclude the range
$\delta=o(k/G)$: there, public outputs with uniform label can already meet
some accuracy requirements, and the exact boundary depends on the internal
states, not only on the winner mass.
